\begin{frame}{3샘플 로지스틱회귀: 단계 1 업데이트와 해석}
  업데이트: $w_0 = 0 - 0.5(-0.1667) = 0.0833$,
  \; $w_1 = 0 - 0.5(-0.3333) = 0.1667$.
  \vskip0.6em
  이제 예측이 더 이상 0.5가 아니다:
  \[
  h(1) = \sigma(0.0833 + 0.1667) = 0.5619, \hspace{1em}
  h(2) = \sigma(0.0833 + 0.3333) = 0.5381, \hspace{1em}
  h(3) = \sigma(0.0833 + 0.5) = 0.6418
  \]
  손실은 \textbf{0.6475} 로 떨어짐.
  \vskip0.6em
  \begin{block}{해석}
    두 샘플 ($x = 1, 3$) 은 양성이므로 모두 $h > 0.5$ 로 ``정답
    쪽''. 유일한 음성 ($x = 2$) 은 가장 작은 기울기로 0.5를 겨우
    넘겨 예측.
    \\ \kb{``반쯤 맞고 반쯤 틀린 예측''을 받으면 그래디언트도
    그만큼 작아진다} --- 이것이 직관.
  \end{block}
  {\footnotesize 실습 노트북 §2 (numpy, 300 epoch):
  $\boldsymbol{w} = [-1.494, 1.345, 0.608]$, 최종 손실 0.3626 ---
  sklearn \texttt{LogisticRegression(C=1e6)} 와 거의 같은
  $[-1.573, 1.492, 0.662]$ 확인.}
\end{frame}
